% GENERATED FILE. Edit canonical lessons, then run npm run generate:books.
% canonical-source-hash: fnv1a64:b5a521202f27c9ce
% canonical-lessons: ML-C529-vivekam, ML-C529-tanretam, ML-C529-niskalankamaya, ML-C529-dussilam, ML-C529-sahisnuta

\chapter{Good Sense and Bad Habits}
\label{ch:ml-good-sense-and-bad-habits}
\hlchaptermodality{529}

\begin{chapteropening}
\textbf{By the end of this chapter:} \emph{I can say good sense, self-assurance, guileless, a bad habit and tolerance.}

Complete the last lesson of chapter 529: I can say good sense, self-assurance, guileless, a bad habit and tolerance.
\end{chapteropening}

\section[\texorpdfstring{\ml{വിവേകം} (vivēkaṁ)}{വിവേകം (vivēkaṁ)}]{\ml{വിവേകം} (vivēkaṁ) — good sense; wisdom}

\label{lesson:ML-C529-vivekam}



\begin{quote}
Before the new one: say the Malayalam for sincerity, then the Malayalam for discipline.
\end{quote}

\subsection*{You'll want to know: \ml{വിവേകം}}
\textbf{\ml{വിവേകം}} — \emph{vivēkaṁ} — \textquotedblleft{}good sense; wisdom\textquotedblright{}.

\begin{quote}
\ml{പ്രശ്നം} \ml{വന്നപ്പോൾ} \ml{അവൾ} \ml{വിവേകത്തോടെ} \ml{തീരുമാനമെടുത്തു}.
\end{quote}

\begin{quote}
\emph{praśnaṁ vannappōḷ avaḷ vivēkattōṭe tīrumānameṭuttu.}
\end{quote}

\begin{quote}
When the problem came, she decided with good sense.
\end{quote}

\ml{വിവേകമുള്ള} means sensible; \ml{വിവേകമില്ലാത്ത}, foolish.

\begin{cousinweb}
\textbf{\ml{വിവേകം}} (\emph{vivēkaṁ}) is borrowed from Sanskrit \emph{viveka}, \textquotedblleft{}discrimination, telling things apart\textquotedblright{} (Monier-Williams): good sense is knowing right from wrong.
\end{cousinweb}

\begin{grammarlens}[title={What you've built}]
One more word for describing what you lived and felt.
\end{grammarlens}

\subsection*{Your turn}
\begin{itemize}
\raggedright
  \item \emph{Say it:} \emph{vivēkaṁ}
  \item \emph{Say it:} \emph{vivēkaṁ}, once more
  \item \emph{Say it:} \emph{praśnaṁ vannappōḷ avaḷ vivēkattōṭe tīrumānameṭuttu.} — When the problem came, she decided with good sense.
  \item \emph{Recall:} say the Malayalam for sincerity, then the Malayalam for discipline, then say \emph{vivēkaṁ} again
\end{itemize}

\begin{tcolorbox}[breakable,colback=teal!4,colframe=teal!35!black,title={Before you move on}]
What does \ml{വിവേകം} mean? (\textquotedblleft{}Good sense; wisdom\textquotedblright{}.) Say it once more.
\end{tcolorbox}
\section[\texorpdfstring{\ml{തന്റേടം} (tanṟēṭaṁ)}{തന്റേടം (tanṟēṭaṁ)}]{\ml{തന്റേടം} (tanṟēṭaṁ) — self-assurance; boldness}

\label{lesson:ML-C529-tanretam}



\begin{quote}
Before the new one: say the Malayalam for discipline, then the Malayalam for good sense.
\end{quote}

\subsection*{You'll want to know: \ml{തന്റേടം}}
\textbf{\ml{തന്റേടം}} — \emph{tanṟēṭaṁ} — \textquotedblleft{}self-assurance; boldness\textquotedblright{}.

\begin{quote}
\ml{ഒറ്റയ്ക്ക്} \ml{വിദേശത്ത്} \ml{പോകാൻ} \ml{അവൾക്ക്} \ml{തന്റേടം} \ml{ഉണ്ട്}.
\end{quote}

\begin{quote}
\emph{oṟṟaykkŭ vidēśattŭ pōkān avaḷkkŭ tanṟēṭaṁ uṇṭŭ.}
\end{quote}

\begin{quote}
She has the nerve to go abroad on her own.
\end{quote}

Praise for a confident person; \ml{തന്റേടി} can also mean a cheeky one.

\begin{cousinweb}
\textbf{\ml{തന്റേടം}} (\emph{tanṟēṭaṁ}) seems to begin with \ml{തൻ}, \textquotedblleft{}self\textquotedblright{}, a native word (DEDR 3196); the rest of the word is not settled.
\end{cousinweb}

\begin{grammarlens}[title={What you've built}]
One more word for describing what you lived and felt.
\end{grammarlens}

\subsection*{Your turn}
\begin{itemize}
\raggedright
  \item \emph{Say it:} \emph{tanṟēṭaṁ}
  \item \emph{Say it:} \emph{tanṟēṭaṁ}, once more
  \item \emph{Say it:} \emph{oṟṟaykkŭ vidēśattŭ pōkān avaḷkkŭ tanṟēṭaṁ uṇṭŭ.} — She has the nerve to go abroad on her own.
  \item \emph{Recall:} say the Malayalam for discipline, then the Malayalam for good sense, then say \emph{tanṟēṭaṁ} again
\end{itemize}

\begin{tcolorbox}[breakable,colback=teal!4,colframe=teal!35!black,title={Before you move on}]
What does \ml{തന്റേടം} mean? (\textquotedblleft{}Self-assurance; boldness\textquotedblright{}.) Say it once more.
\end{tcolorbox}
\section[\texorpdfstring{\ml{നിഷ്കളങ്കമായ} (niṣkaḷaṅkamāya)}{നിഷ്കളങ്കമായ (niṣkaḷaṅkamāya)}]{\ml{നിഷ്കളങ്കമായ} (niṣkaḷaṅkamāya) — guileless; innocent}

\label{lesson:ML-C529-niskalankamaya}



\begin{quote}
Before the new one: say the Malayalam for good sense, then the Malayalam for self-assurance.
\end{quote}

\subsection*{You'll want to know: \ml{നിഷ്കളങ്കമായ}}
\textbf{\ml{നിഷ്കളങ്കമായ}} — \emph{niṣkaḷaṅkamāya} — \textquotedblleft{}guileless; innocent\textquotedblright{}.

\begin{quote}
\ml{കുഞ്ഞിന്റെ} \ml{നിഷ്കളങ്കമായ} \ml{ചിരി} \ml{എല്ലാവർക്കും} \ml{ഇഷ്ടപ്പെട്ടു}.
\end{quote}

\begin{quote}
\emph{kuññinṟe niṣkaḷaṅkamāya ciri ellāvarkkuṁ iṣṭappeṭṭu.}
\end{quote}

\begin{quote}
Everyone loved the baby's innocent smile.
\end{quote}

\ml{നിഷ്}- without + \ml{കളങ്കം} a stain: without a stain, guileless.

\begin{cousinweb}
\textbf{\ml{നിഷ്കളങ്കമായ}} (\emph{niṣkaḷaṅkamāya}) is built from Sanskrit \emph{niṣkalaṅka}, \textquotedblleft{}stainless, spotless\textquotedblright{} (Monier-Williams), with the Malayalam adjective ending -\ml{ആയ}.
\end{cousinweb}

\begin{grammarlens}[title={What you've built}]
One more word for describing what you lived and felt.
\end{grammarlens}

\subsection*{Your turn}
\begin{itemize}
\raggedright
  \item \emph{Say it:} \emph{niṣkaḷaṅkamāya}
  \item \emph{Say it:} \emph{niṣkaḷaṅkamāya}, once more
  \item \emph{Say it:} \emph{kuññinṟe niṣkaḷaṅkamāya ciri ellāvarkkuṁ iṣṭappeṭṭu.} — Everyone loved the baby's innocent smile.
  \item \emph{Recall:} say the Malayalam for good sense, then the Malayalam for self-assurance, then say \emph{niṣkaḷaṅkamāya} again
\end{itemize}

\begin{tcolorbox}[breakable,colback=teal!4,colframe=teal!35!black,title={Before you move on}]
What does \ml{നിഷ്കളങ്കമായ} mean? (\textquotedblleft{}Guileless; innocent\textquotedblright{}.) Say it once more.
\end{tcolorbox}
\section[\texorpdfstring{\ml{ദുശ്ശീലം} (duśśīlaṁ)}{ദുശ്ശീലം (duśśīlaṁ)}]{\ml{ദുശ്ശീലം} (duśśīlaṁ) — a bad habit}

\label{lesson:ML-C529-dussilam}



\begin{quote}
Before the new one: say the Malayalam for self-assurance, then the Malayalam for guileless.
\end{quote}

\subsection*{You'll want to know: \ml{ദുശ്ശീലം}}
\textbf{\ml{ദുശ്ശീലം}} — \emph{duśśīlaṁ} — \textquotedblleft{}a bad habit\textquotedblright{}.

\begin{quote}
\ml{പുകവലി} \ml{ഒരു} \ml{ദുശ്ശീലമാണ്}.
\end{quote}

\begin{quote}
\emph{pukavali oru duśśīlamāṇŭ.}
\end{quote}

\begin{quote}
Smoking is a bad habit.
\end{quote}

\ml{ദു}- bad + \ml{ശീലം}, a habit, the word already familiar.

\begin{cousinweb}
\textbf{\ml{ദുശ്ശീലം}} (\emph{duśśīlaṁ}) is borrowed from Sanskrit \emph{duḥśīla}, \textquotedblleft{}ill-behaved\textquotedblright{} (CDIAL 6384), built from \emph{duṣ-}, \textquotedblleft{}bad\textquotedblright{}, and \emph{śīla}, \textquotedblleft{}habit, nature\textquotedblright{}.
\end{cousinweb}

\begin{grammarlens}[title={What you've built}]
One more word for describing what you lived and felt.
\end{grammarlens}

\subsection*{Your turn}
\begin{itemize}
\raggedright
  \item \emph{Say it:} \emph{duśśīlaṁ}
  \item \emph{Say it:} \emph{duśśīlaṁ}, once more
  \item \emph{Say it:} \emph{pukavali oru duśśīlamāṇŭ.} — Smoking is a bad habit.
  \item \emph{Recall:} say the Malayalam for self-assurance, then the Malayalam for guileless, then say \emph{duśśīlaṁ} again
\end{itemize}

\begin{tcolorbox}[breakable,colback=teal!4,colframe=teal!35!black,title={Before you move on}]
What does \ml{ദുശ്ശീലം} mean? (\textquotedblleft{}A bad habit\textquotedblright{}.) Say it once more.
\end{tcolorbox}
\section[\texorpdfstring{\ml{സഹിഷ്ണുത} (sahiṣṇuta)}{സഹിഷ്ണുത (sahiṣṇuta)}]{\ml{സഹിഷ്ണുത} (sahiṣṇuta) — tolerance}

\label{lesson:ML-C529-sahisnuta}



\begin{quote}
Before the new one: say the Malayalam for guileless, then the Malayalam for a bad habit.
\end{quote}

\subsection*{You'll want to know: \ml{സഹിഷ്ണുത}}
\textbf{\ml{സഹിഷ്ണുത}} — \emph{sahiṣṇuta} — \textquotedblleft{}tolerance\textquotedblright{}.

\begin{quote}
\ml{മറ്റുള്ളവരുടെ} \ml{അഭിപ്രായത്തോട്} \ml{സഹിഷ്ണുത} \ml{വേണം}.
\end{quote}

\begin{quote}
\emph{maṟṟuḷḷavaruṭe abhiprāyattōṭŭ sahiṣṇuta vēṇaṁ.}
\end{quote}

\begin{quote}
One needs tolerance towards other people's opinions.
\end{quote}

From \ml{സഹിഷ്ണു}, able to bear, the same root as \ml{സഹിക്കുക}.

\begin{cousinweb}
\textbf{\ml{സഹിഷ്ണുത}} (\emph{sahiṣṇuta}) is borrowed from Sanskrit \emph{sahiṣṇutā}, the noun of \emph{sahiṣṇu}, \textquotedblleft{}patient, enduring\textquotedblright{} (Monier-Williams), from the root \emph{sah}, \textquotedblleft{}to bear\textquotedblright{}.
\end{cousinweb}

\begin{grammarlens}[title={What you've built}]
One more word for describing what you lived and felt.
\end{grammarlens}

\subsection*{Your turn}
\begin{itemize}
\raggedright
  \item \emph{Say it:} \emph{sahiṣṇuta}
  \item \emph{Say it:} \emph{sahiṣṇuta}, once more
  \item \emph{Say it:} \emph{maṟṟuḷḷavaruṭe abhiprāyattōṭŭ sahiṣṇuta vēṇaṁ.} — One needs tolerance towards other people's opinions.
  \item \emph{Recall:} say the Malayalam for guileless, then the Malayalam for a bad habit, then say \emph{sahiṣṇuta} again
\end{itemize}

\begin{tcolorbox}[breakable,colback=teal!4,colframe=teal!35!black,title={Before you move on}]
What does \ml{സഹിഷ്ണുത} mean? (\textquotedblleft{}Tolerance\textquotedblright{}.) Say it once more.
\end{tcolorbox}
